% !TeX root = ../../poisson-del-pezzo.tex
\section{零 Poisson 構造の特異性と標準的な解消}
\label{sec:singularities}

次数2の相対 Jacobi 錐には，非零部分の滑らかさより強い幾何学的構造がある．
本節では，零 Poisson 構造を保つ座標変換に関して標準的な blow-up を構成し，
モジュライスタックに降下させる．以下では $n\geq2$，$m=n+1$ とし，
$C_m=\operatorname{Cone}(\Gr(2,m))\subset\bigwedge^2\C^m$ と書く．
Grassmann 多様体の Pl\"ucker 埋め込みが射影正規で arithmetically
Cohen--Macaulay であることは標準単項式理論の既知の結果として使う
\cite{Svanes1973,DeConciniProcesi1981}．

\begin{lemma}[Pl\"ucker 錐の局所完全交叉性]
\label{lemma:cone-lci}
$C_m$ の頂点 $0$ において
\[
 \operatorname{edim}_0 C_m=\binom m2,\qquad
 \dim_0 C_m=2m-3,\qquad
 \operatorname{codim}_0 C_m=\frac{(m-2)(m-3)}2.
\]
頂点の定義イデアルの最小生成元数は $\binom m4$ である．
従って $m=3$ では $C_m=\mathbb A^3$，$m=4$ では非退化な二次超曲面の錐，
$m\geq5$ では $C_m$ は頂点で局所完全交叉でない．
\end{lemma}
\begin{proof}
錐の座標環を $R=\C[a_{ij}]/I$ とする．$I$ は斉次二次の Pl\"ucker 式
\[
 q_{ijkl}=a_{ij}a_{kl}-a_{ik}a_{jl}+a_{il}a_{jk}
 \qquad (i<j<k<l)
\]
で生成される．これらは相異なる四つの添字集合ごとに，互いに共有しない
二次単項式を含むので一次独立であり，$I/\mathfrak m I$ の基底をなす．
ゆえに最小生成元数は $\binom m4$ である．
$I$ に一次式がないので，接空間次元は $\binom m2$ である．
$\dim\Gr(2,m)=2(m-2)$ から錐の次元は $2m-3$ となる．
$m=3$ の場合は $I=0$，$m=4$ では唯一の Pl\"ucker 式は
六変数の非退化二次形式である．$m\geq5$ の場合には
\[
 \binom m4>\binom m2-(2m-3)
 =\frac{(m-2)(m-3)}2
\]
が成り立つので，正則局所環における完全交叉の生成元数の必要条件に反する．
\end{proof}

\begin{lemma}[頂点の blow-up]
\label{lemma:cone-blowup}
$Y_m=\Gr(2,m)$，$L=\mathcal O_{Y_m}(-1)$ とおく．
$C_m$ の頂点の blow-up は，標準的に
\[
 \operatorname{Bl}_0(C_m)\simeq\operatorname{Tot}_{Y_m}(L)
\]
である．例外因子は零切断 $Y_m$ であり，その法線束は $L$ である．
この同型は $\operatorname{GL}_m$ および錐のスカラー作用に関して同変である．
\end{lemma}
\begin{proof}
$C_m\setminus\{0\}$ の点は，Pl\"ucker 直線 $[v]\in Y_m$ と
その直線の非零ベクトル $v$ の組に一意に対応する．
従って閉部分スキーム
\[
 \{(v,[\ell])\in C_m\times Y_m\mid v\in\ell\}
\]
は $\operatorname{Tot}(L)$ である．一方，頂点のイデアルは一次座標
$a_{ij}$ 全体で生成され，その blow-up は $C_m$ 上のこれらの座標による
有理写像 $C_m\dashrightarrow\mathbb P(\bigwedge^2\C^m)$ のグラフの閉包である．
$C_m\setminus\{0\}$ ではこの写像の像は $Y_m$ に入り，閉包を取っても
$Y_m$ の斉次 Pl\"ucker 方程式が保たれる．上の閉部分スキームは各
Pl\"ucker チャート上で「一つの座標が自由なスカラー，残りはその座標に
Grassmann チャート座標を掛けたもの」という形だから，グラフ閉包に等しい．
すべての構成が線型作用に関して自然である．
\end{proof}

\begin{theorem}[相対零切断の blow-up による解消]
\label{thm:poisson-resolution}
次数2のモジュライの零切断を
$s:\mathcal D_{n,2}\hookrightarrow\mathcal P_{n,2}$ とおく．
$\mathcal R_{n,2}=\operatorname{Bl}_{s(\mathcal D_{n,2})}\mathcal P_{n,2}$
とすると，次が成り立つ．
\begin{enumerate}
\item $\mathcal R_{n,2}$ は有限型の滑らかな Deligne--Mumford スタックであり，
$\mathcal R_{n,2}\to\mathcal P_{n,2}$ は representable，射影的かつ固有で，
零切断の補集合では同型である．$n\geq3$ のときこれは特異性の解消を与える．
\item 例外因子 $\mathcal E_{n,2}$ は $\mathcal D_{n,2}$ 上の滑らかな相対
Grassmann 族
\[
 \mathcal E_{n,2}\simeq
 [\mathcal P_{n,2}^{\times}/\mathbb G_m]
\longrightarrow\mathcal D_{n,2}
\]
に標準的に同型である．相対 Poisson パラメータのベクトル束を
$\mathcal V$，その射影 Grassmann 部分を $\mathcal G\subset\mathbb P(\mathcal V)$
と書けば，解消は $\mathcal G$ 上の tautological line bundle
$\mathcal O_{\mathcal G}(-1)$ の全空間に一致する．
例外因子の法線束は $\mathcal O_{\mathcal G}(-1)$ である．
\item 任意の $S\to\mathcal D_{n,2}$ に対して，この解消は対応する
相対 Pl\"ucker 錐の零切断の blow-up と可換である．
\end{enumerate}
\end{theorem}
\begin{proof}
命題~\ref{prop:intrinsic-bivector-bundle} と
系~\ref{cor:projective-moduli} により，$\mathcal V$ と $\mathcal G$ は
偏極線束を基底から捻る操作に対して不変であり，$\mathcal D_{n,2}$ 上に降下する．
$\mathcal D_{n,2}$ の標準方程式 atlas $U_2$ 上では，
$\mathcal P_{n,2}$ は $U_2\times C_m$，零切断は $U_2\times\{0\}$ である．
blow-up は平坦基底変換と可換なので，補題~\ref{lemma:cone-blowup} から
\[
 \mathcal R_{n,2}\times_{\mathcal D_{n,2}}U_2
 \simeq U_2\times\operatorname{Tot}_{Y_m}\mathcal O_{Y_m}(-1)
\]
を得る．同型は式~\eqref{eq:moduli-action-two} の $G_2$ の作用と同変である．
従って商スタックへ降下し，滑らかさ，例外因子と法線束の記述を得る．
blow-up の射影性・固有性と零切断外での同型性も atlas 上で確認して降下する．
記載した基底変換は各 atlas 上での自明な積表示と blow-up の構成から従う．
\end{proof}

\begin{proposition}[Pl\"ucker 錐の標準特異性]
\label{prop:plucker-terminal}
$m\geq4$ のとき $C_m$ は頂点のみで特異であり，その特異点は
Gorenstein terminal（従って rational）である．
補題~\ref{lemma:cone-blowup} の解消 $\rho:\widetilde C_m\to C_m$ の例外因子を
$E$ とすると
\[
 K_{\widetilde C_m}=\rho^*K_{C_m}+(m-1)E
\]
という discrepancy 等式が成立する．特に $m=n+1$ では例外因子の
discrepancy は $n$ である．これは Poisson モジュライの滑らかな
標準方程式 atlas に関する主張であって，coarse moduli space の
商特異点についての主張ではない．
\end{proposition}
\begin{proof}
Grassmann 多様体の普遍部分束 $\mathcal S$ と商束 $\mathcal Q$ から
$T_{Y_m}=\mathcal S^*\otimes\mathcal Q$，従って
$\omega_{Y_m}\simeq\mathcal O_{Y_m}(-m)$ である．
Pl\"ucker 座標環 $R$ は射影正規かつ Cohen--Macaulay であり，
その次数付き canonical module は
\[
 \omega_R\simeq
 \bigoplus_{k\in\mathbb Z}H^0(Y_m,\omega_{Y_m}(k))
 \simeq R(-m)
\]
である．ここで最後の等式では $\omega_{Y_m}=\mathcal O(-m)$ と
Pl\"ucker 線束の切断環の記述を用いた．従って $C_m$ は Gorenstein であり，
$K_{C_m}$ は（非斉次な意味で）自明である．

$\widetilde C_m=\operatorname{Tot}_{Y_m}(L)$ に対して
$\omega_{\widetilde C_m}\simeq
p^*(\omega_{Y_m}\otimes L^{-1})=p^*\mathcal O_{Y_m}(-(m-1))$ である．
零切断 $E$ の法線束が $L=\mathcal O(-1)$ なので
$\mathcal O_{\widetilde C_m}(E)=p^*\mathcal O(-1)$ であり，
discrepancy は $m-1$ である．頂点以外の錐は滑らかだから，
他の例外因子についても滑らかな解消の上でさらに吹き上げれば
discrepancy が正になる．従って terminal である．
標数0で terminal は klt，klt は rational である
\cite[Chapter 5]{KollarMori1998}．
\end{proof}

\begin{corollary}[相対モジュライの特異性の次元別分類]
\label{cor:moduli-singularity-types}
次数2の $\mathcal P_{n,2}\to\mathcal D_{n,2}$ は，
$n=2$ では滑らか，$n=3$ では零切断に沿い相対的に通常の
五次元二次錐特異点を持つ hypersurface 型，
$n\geq4$ では零切断に沿い非局所完全交叉である．
$n\geq3$ の零点における相対 Zariski 接空間の次元超過は
\[
 \binom{n+1}{2}-(2n-1)=\frac{(n-1)(n-2)}2
\]
である．どの $n\geq3$ でも滑らかな atlas の特異性は terminal であり，
その零切断 blow-up は定理~\ref{thm:poisson-resolution} の滑らかな解消となる．
\end{corollary}
\begin{proof}
相対分類，補題~\ref{lemma:cone-lci}，
定理~\ref{thm:poisson-resolution}，命題~\ref{prop:plucker-terminal} を合わせる．
\end{proof}

\begin{remark}[既知の Grassmann 幾何と本稿の寄与]
錐の blow-up と Gorenstein 性，標準単項式理論に基づく局所環の性質は
既知の代数幾何である．ここでの寄与は，次数2の相対 Jacobi スキームとの
スキーム論的同定を経由し，零 Poisson 切断とその解消を
$G_2$-同変に具体化して，Poisson del Pezzo のモジュライスタックに
降下させた点にある．この降下による新規性の優先権は未確認である．
\end{remark}
